% !TeX root = ../教材题库_全部条目_转写工作稿.tex
% Source: XJ-B1-C03-TH-P083-Y0277
% 第3章 / 3.2 函数的基本性质 / 教材 p.83 / PDF 第90页
% OCR未校对：公式与文字均可能有误；未完成数学TeX转写；禁止直接发给学生
% 原题及相邻原解截图（核对用）：../题目原图/XJ-B1-C03-TH-P083-Y0277-01.webp
% 本文件只在转写完成后作为题干使用；OCR 条目尚不是可靠题干。
\begin{bankitem}{p.83 思考（含紧邻说明）}{第3章 / 3.2 函数的基本性质 / 教材 p.83 / PDF 第90页}{XJ-B1-C03-TH-P083-Y0277}{OCR未校对：公式与文字均可能有误；未完成数学TeX转写；禁止直接发给学生}
\begin{ocrdraft}
[教材p.83 / PDF第90页]
\par
？思考
\par
（D观察图3．2一5中的两个函数图象，它们有什么共同特征？
\par
（2）观察图3，2一6中的两个函数图象，它们有什么共同特征？
\par
y=-2x
\par
Y=
\par
0
\par
0
\par
y：2、
\par
图3．2一5
\par
图3．2一6
\par
不难发现，图3．2一5中的两个图象，都是以3，轴为对称轴的轴对称图形；图
\par
3．2一6中的两个图象，都是以原点为中心的中心对称图形．
\par
如果F(x）的图象是以3，轴为对称轴的轴对称图形，就称F(x）是偶函数．如果
\par
F(x）的图象是以原点为中心的中心对称图形，就称F(x）是奇函数．
\end{ocrdraft}
\end{bankitem}
